\begin{abstract}
The \textit{Shortest Path Problem (SPP)} is a fundamental problem in graph theory and computer science, with many applications (e.g. navigation, network routing).
This Thesis involves the verification of the \textit{DIJKSTRA-PREDICTION} algorithm, 
which utilizes machine learning techniques that was proposed by \textit{Willem Feijen and Guido Schäfer} for the substantial reduction of priority queue operations in regards to \textit{Dijkstra's algorithm}, 
specifically targeting the \textit{Single-Source Many-Targets Shortest-Path Problem (SSMTSP)}. Further more the algorithm proposed can reduce running time in many scenarios. 
This thesis aims to further contribute in the Machine Learning approaches for the design of improved algorithms such as those that are known as \textit{Algorithms with Predictions} (or \textit{Learning Augmented Algorithms}). 
We include a theoretical analysis of the algorithm, several experiments as well as an interactive visualization tool to further validate the findings which allows for the free design of a graph (or random generation) 
as well as a "trap" scenario to further emphasize on the efficiency of the algorithm compared to classical methods.
\end{abstract}